% !TEX root = main.tex
\chapter{COMPLEX ANALYSIS, RATIONAL AND MEROMORPHIC ASYMPTOTICS}
\label{chap:complex analysis, rational and meromorphic asymptotics}



\section{Generating functions as analytic objects}
% Note IV.6.
\vspace{\baselineskip}
\noindent
$\rhd$ \textbf{IV.6.} \textit{Integrals of rational fractions.}

\vspace{\baselineskip}
\noindent
\begin{equation*}
    J_{m} = \int_{-\infty}^{\infty}\frac{dx}{(1 + x^{2})^{m}}
\end{equation*}
Define $f(z) = \frac{1}{(1 + z^{2})^{m}} = \frac{1}{(z - i)^{m}(z + i)^{m}}$. We integrate over the contour $\Gamma_{R}$ consisting of the real
interval $[-R, R]$ and the upper semicircle $C_{R} = \{ z = Re^{i\theta}: 0 \leq \theta \leq \pi \}$.

\vspace{\baselineskip}
\noindent
For $z \in C_{R}$ and $R > 1$, $|z^{2} + 1| \geq |z^{2}| - 1 = R^{2} - 1$. Thus
\begin{equation*}
    |f(z)| \leq \frac{1}{(R^{2} - 1)^{m}}
\end{equation*}
The length of $C_{R}$ is $\pi R$, so
\begin{equation*}
    \left|\int_{C_{R}}f(z)dz\right| \leq \frac{\pi R}{(R^{2} - 1)^{m}} \to 0 \; \text{as} \; R \to \infty
\end{equation*}
Hence
\begin{equation*}
    \lim_{R \to \infty}\int_{\Gamma_{R}}f(z)dz = \int_{-\infty}^{\infty}\frac{dx}{(1 + x^{2})^{m}} = J_{m}
\end{equation*}
For $R > 1$, the only singularity inside $\Gamma_{R}$ is the pole of order $m$ at $z = i$. By the residue theorem,
\begin{equation*}
    \int_{\Gamma_{R}}f(z)dz = 2\pi i\text{Res}(f; i)
\end{equation*}
Compute the residue at $z = i$:
\begin{equation*}
    \text{Res}(f; i) = \frac{1}{(m - 1)!}\lim_{z \to i}\frac{d^{m - 1}}{dz^{m - 1}}[(z - i)^{m}f(z)]
        = \frac{1}{(m - 1)!}\lim_{z \to i}\frac{d^{m - 1}}{dz^{m - 1}}(z + i)^{-m}
\end{equation*}
The $k$-th derivative of $g(z) \coloneqq (z + i)^{-m}$ is
\begin{equation*}
    g^{(k)}(z) = (-1)^{k}\frac{(m + k - 1)!}{(m - 1)!}(z + i)^{-m - k}
\end{equation*}
For $k = m - 1$:
\begin{equation*}
    g^{(m - 1)}(z) = (-1)^{m - 1}\frac{(2m - 2)!}{(m - 1)!}(z + i)^{-2m + 1}
\end{equation*}
Evaluating at $z = i$ gives
\begin{equation*}
    g^{(m - 1)}(i) = (-1)^{m - 1}\frac{(2m - 2)!}{(m - 1)!}(2i)^{-2m + 1}
\end{equation*}
Hence
\begin{equation*}
    \begin{aligned}
        \text{Res}(f; i)
        &= \frac{1}{(m - 1)!}(-1)^{m - 1}\frac{(2m - 2)!}{(m - 1)!}(2i)^{-2m + 1} \\
        &= (-1)^{m - 1}\frac{(2m - 2)!}{[(m - 1)!]^{2}}\frac{(-1)^{m}i}{2^{2m - 1}} \\
        &= -i\frac{(2m - 2)!}{[(m - 1)!]^{2}2^{2m - 1}} \\
    \end{aligned}
\end{equation*}
\begin{equation*}
    J_{m} = 2\pi i\left(-i\frac{(2m - 2)!}{[(m - 1)!]^{2}2^{2m - 1}}\right) = \frac{\pi}{2^{2m - 2}}\binom{2m - 2}{m - 1}
        = \frac{\pi}{4^{m - 1}}\binom{2m - 2}{m - 1}
\end{equation*}

\vspace{2\baselineskip}
\noindent
\begin{equation*}
    K_{m} = \int_{-\infty}^{\infty}\frac{dx}{(1^{2} + x^{2})(2^{2} + x^{2}) \cdots (m^{2} + x^{2})}
\end{equation*}
Define $f(z) = \frac{1}{\prod_{k = 1}^{m}(z^{2} + k^{2})}$. We integrate over the contour $\Gamma_{R}$ consisting of the real
interval $[-R, R]$ and the upper semicircle $C_{R} = \{ z = Re^{i\theta}: 0 \leq \theta \leq \pi \}$.

\vspace{\baselineskip}
\noindent
For $R > m$, the poles inside $\Gamma_{R}$ are $z_{k} = ik, k = 1, \cdots m$. On $C_{R}$, $|z^{2} + k^{2}| \geq R^{2} - k^{2} \geq (R/2)^{2}$ eventually,
so $|f(z)| \leq \text{const}/R^{2m}$ and the length of $C_{R}$ is $\pi R$. Hence
\begin{equation*}
    \lim_{R \to \infty}\int_{C_{R}}f(z)dz = 0
\end{equation*}
and by the residue theorem
\begin{equation*}
    K_{m} = 2\pi i\sum_{k = 1}^{m}\text{Res}(f; ik)
\end{equation*}
Since the poles are simple
\begin{equation*}
    \text{Res}(f; ik) = \lim_{z \to ik}(z - ik)f(z) = \frac{1}{2ik}\prod_{j \neq k}\frac{1}{(ik)^{2} + j^{2}}
        = \frac{1}{2ik}\prod_{j \neq k}\frac{1}{j^{2} - k^{2}}
\end{equation*}
Thus
\begin{equation*}
    K_{m} = 2\pi i\sum_{k = 1}^{m}\frac{1}{2ik}\prod_{j \neq k}\frac{1}{j^{2} - k^{2}} 
        = \pi\sum_{k = 1}^{m}\frac{1}{k}\prod_{j \neq k}\frac{1}{j^{2} - k^{2}}
\end{equation*}
Factor the product:
\begin{equation*}
    \prod_{j \neq k}j^{2} - k^{2} = \prod_{j \neq k}(j - k)\prod_{j \neq k}(j + k)
\end{equation*}
For fixed $k$,
\begin{equation*}
    \prod_{j \neq k}(j - k) = (1 - k)(2 - k) \cdots ((k - 1) - k)(1)(2) \cdots (m - k) = (-1)^{k - 1}(k - 1)!(m - k)!
\end{equation*}
\begin{equation*}
    \prod_{j \neq k}(j + k) = (1 + k)(2 + k) \cdots ((k - 1) + k)(2k + 1)(2k + 2) \cdots (m + k) = \frac{(m + k)!}{k!(2k)}
\end{equation*}
Therefore
\begin{equation*}
    \frac{1}{k}\prod_{j \neq k}\frac{1}{j^{2} - k^{2}} = \frac{1}{k}\left((-1)^{k - 1}(k - 1)!(m - k)! \cdot \frac{(m + k)!}{k!(2k)}\right)^{-1}
        = \frac{(-1)^{k - 1}2k}{(m - k)!(m + k)!}
\end{equation*}
and
\begin{equation*}
    K_{m} = \pi\sum_{k = 1}^{m}\frac{(-1)^{k - 1}2k}{(m - k)!(m + k)!} = \frac{2\pi}{(2m)!}\sum_{k = 1}^{m}(-1)^{k - 1}k\binom{2m}{m - k}
\end{equation*}

\vspace{\baselineskip}
\noindent
To simplify the sum $S \coloneqq \sum_{k = 1}^{m}(-1)^{k - 1}k\binom{2m}{m + k}$, we can use the method of generating functions (what else?).
Recall the coefficient of $z^{m}$ in the product of two polynomials $A(z) = \sum a_{k}z^{k}$ and $B(z) = \sum b_{j}z^{j}$ is given by
\begin{equation*}
    [z^{m}]A(z)B(z) = \sum_{k}a_{k}b_{m - k}
\end{equation*}
Let's choose $A(z) = \sum_{k = 1}^{m}(-1)^{k - 1}kz^{k}$ and $B(z) = (1 + z)^{2m} = \sum_{j = 0}^{2m}\binom{2m}{j}z^{j}$. Then, the coefficient
of $z^{m}$ in the product $A(z)B(z)$ is precisely our sum $S$:
\begin{equation*}
    S = [z^{m}]\sum_{k = 1}^{m}(-1)^{k - 1}kz^{k}(1 + z)^{2m}
\end{equation*}
Since we only care about the coefficient of $z^{m}$, any terms in $A(z)$ with powers of $z$ strictly greater than $m$ will not contribute to the $z^{m}$ term
when multiplied by $B(z)$ (as $B(z)$ contains only non-negative powers of $z$). Therefore, we can extend the upper limit of the sum in $A(z)$ to infinity
without changing the coefficient of $z^{m}$:
\begin{equation*}
    A(z) \equiv \sum_{k = 1}^{\infty}(-1)^{k - 1}kz^{k} = \frac{z}{(1 + z)^{2}} \mod z^{m + 1}
\end{equation*}
Hence
\begin{equation*}
    S = [z^{m}]\frac{z}{(1 + z)^{2}} \cdot (1 + z)^{2m} = [z^{m}]z(1 + z)^{2m - 2} = \binom{2m - 2}{m - 1}
\end{equation*}

\vspace{\baselineskip}
\noindent
Finally
\begin{equation*}
    K_{m} = \frac{2\pi}{(2m)!}\binom{2m - 2}{m - 1} = \frac{\pi}{m(2m - 1)[(m - 1)!]^{2}}
\end{equation*}



\section{Analytic functions and meromorphic functions}



\section{Singularities and exponential growth of coefficients}



\section{Closure properties and computable bounds}



\section{Rational and meromorphic functions}



\section{Localization of singularities}



\section{Singularities and functional equations}



\section{Perspective}